\section{Dimensions of comparative analysis}
\label{app:dimensions}

The patterns in this article fall along several dimensions that should remain
separate in analysis. Collapsing them into a single score would hide the
difference between factual error, generic style, and a mismatch with one
reader's needs.

\subsection{Semantic fidelity}

The first dimension concerns preservation of meaning. A revision is faithful
when its factual claims and examples remain grounded in the source, including
statistics, quotations, citations, and technical distinctions. Calibrated
uncertainty belongs to this dimension because a change in evidential verb can
alter the claim even when every noun and number remains the same.

Verb semantics provide a related test. The subject must be able to perform the
named action, literally or within an established technical convention, and its
object or complement must fit the relation expressed by the verb. Evidence can
support a claim; stated conditions may limit its extent or applicability. If
the question concerns implication, the relevant boundary belongs to the
consequences of the claim. Metaphorical movement or agency becomes a defect
when it obscures one of these relations.

\subsection{Reader model and information structure}

Reader-centered prose makes its assumptions about prior knowledge and intended
use recoverable from the selection of detail. At sentence scale, familiar
context tends to appear early enough to connect with the preceding discourse,
while the result or distinction receives an appropriate position of emphasis.
At paragraph scale, evidence and qualifications remain visibly subordinate to
the claim they support.

Several sentence patterns can disrupt this structure. Question-word wrappers
delay the point, while distant subject-verb pairs postpone syntactic
resolution. Unnamed actors and stiff relative clauses create different
problems, so their importance depends on the reading effect rather than the
presence of a particular token.

\subsection{Rhetorical form}

Generic significance, promotional language, vague attribution, decorative
analysis, and optimistic conclusions share a common effect: they ask a
rhetorical form to supply an evaluation that lacks evidential support. Lists
and headings create a structural analogue when they distribute equal weight
across ideas with different roles.

Display text belongs to the same analysis as body prose. Titles, abstracts,
headings, captions, and example labels all frame an interpretation. A generic
title or section-summary abstract can therefore preserve every sentence in
the body while still misrepresenting the argument.

\subsection{Cadence and voice}

Cadence is visible in sentence length, but it also depends on openings, clause
shapes, contrast formulas, and syntactic closure. Several short sentences may
manufacture drama; several medium sentences with identical structures can be
equally monotonous. Repeated formulas such as ``treats X as Y, not Z'' deserve
attention when the syntax persists after the underlying comparison changes.

Voice is not the opposite of regularity. It consists of supported specifics,
recurring judgments, deliberate irregularities, and genre-appropriate
personality. A successful revision reduces mechanical rhythm without
normalizing these choices into the editor's own style.

\subsection{Provenance}

None of the preceding dimensions establishes authorship. Quotations, proper
names, and secondhand examples retain the provenance of their source, while
templates and professional editing can explain a polished surface. Version
history and drafting records provide stronger evidence about production than
stylistic inference. The patterns remain useful because they describe reading
effects even when the author is known.
